% Complete Section 4.1, original PDF 43-44 / printed 33-34.
% Fraktur B/Z are displayed in the user-approved mathcal family; font evidence retained.
% block: hilbert-complexes-duals-heading page:43 type:section
\section{{Hilbert 复形及其对偶}}\label{sec:hilbert-complexes-duals}

% block: hilbert-complexes-subheading page:43 type:subsection
\subsection{{Hilbert 复形}}\label{subsec:hilbert-complexes-definition}

% block: hilbert-complex-definition page:43 type:definition

\begin{Definition}\label{def:hilbert-complex}
{\emph{Hilbert 复形}是一列 Hilbert 空间 $W^k$, 以及一列从 $W^k$ 到 $W^{k+1}$ 的闭的稠密定义线性算子 $d^k$, 满足 $\mathcal R(d^k)\subset\mathcal N(d^{k+1})$.}
\end{Definition}

% block: hilbert-base-domain-spaces page:43 type:ordinary-paragraph
{注意, Hilbert 复形并不仅仅是各空间均为 Hilbert 空间的复形. 在每个层次上, 都有两个 Hilbert 空间: 基空间 $W^k$ 和定义域空间, 后者记为 $V^k$. 由于假设微分 $d^k$ 是闭的, 定义域空间 $V^k$ 配备图范数后就是 Hilbert 空间:}
\[
\|v\|_{V^k}^2=\|v\|_{W^k}^2+\|d^k v\|_{W^{k+1}}^2.
\]

% block: domain-complex-intro page:44 type:ordinary-paragraph
{算子 $d^k$ 可视为从 $V^k$ 到 $V^{k+1}$ 的有界线性映射, 因而它们构成一个有界 Hilbert 复形, 即各空间为 Hilbert 空间且微分为有界线性算子的上链复形, 称为\emph{定义域复形}:}
\begin{equation}\label{eq:hilbert-domain-complex}
0\to V^0\xrightarrow{d}V^1\xrightarrow{d}\cdots\xrightarrow{d}V^n\to0.
\end{equation}

% block: hilbert-versus-domain-complex page:44 type:ordinary-paragraph
{然而, Hilbert 复形 $(W,d)$ 比其定义域复形 $(V,d)$ 包含更多信息. 由前者可构造后者, 但仅知道后者, 无法推知空间 $W^k$ 或 $W^k$ 范数. 除刚才定义的 Hilbert 上链复形外, 当然也可讨论 Hilbert 链复形; 唯一的区别是指标递减.}

% block: closed-fredholm-complexes page:44 type:ordinary-paragraph
{与前面对复形所作的一样, 将微分的像与零空间分别记为 $\mathcal B^k=d^{k-1}V^{k-1}$ 和 $\mathcal Z^k=\mathcal N(d^k)$, 即上边界链空间和闭上链空间. 有 $\mathcal B^k\subset\mathcal Z^k$, 其商为上同调空间 $\mathcal H^k=\mathcal Z^k/\mathcal B^k$. 闭链空间总是闭的, 但边界链空间可能闭, 也可能不闭. 若对所有 $k$, 像空间 $\mathcal B$ 在 $W$ 中闭, 则称 Hilbert 复形是\emph{闭的}. 这当且仅当 $\mathcal B$ 关于 $W$ 范数完备. 由于 $\mathcal Z$ 上的 $V$ 范数与该空间上的 $W$ 范数相同, 等价地可以要求 $\mathcal B$ 关于 $V$ 范数完备, 或要求 $\mathcal B$ 在 $V$ 中闭. 若上同调空间为有限维, 即边界链空间在闭链空间中为有限余维, 则称 Hilbert 复形为 \emph{Fredholm} 复形. 我们遇到的大多数复形都是 Fredholm 的. 由 Theorem~\ref{thm:finite-codimension-closed-range}, Fredholm 复形是闭的.}

% block: reduced-cohomology page:44 type:ordinary-paragraph
{\emergencystretch=2em
{在少数需要处理不闭或尚不知是否闭的 Hilbert 复形的场合, 还引入\emph{约化上同调空间} $\mathcal Z^k/\overline{\mathcal B}^k$, 它是 Hilbert 空间. 当然, 对闭 Hilbert 复形, 它与上同调空间 $\mathcal H^k$ 相同.}

\par}
% block: dual-complex-heading page:44 type:subsection
\subsection{{对偶复形}}\label{subsec:dual-hilbert-complex}

% block: dual-hilbert-complex-domains page:44 type:ordinary-paragraph
{若 $(W,d)$ 为 Hilbert 复形, 则 $d^{k-1}$ 是从 $W^{k-1}$ 到 $W^k$ 的闭的稠密定义算子, 因而其伴随 (记为 $d_k^*$) 是从 $W^k$ 到 $W^{k-1}$ 的闭的稠密定义算子. 记其定义域为 $V_k^*$, 它是 $W^k$ 的稠密子集, 一般与 $V^k$ 很不相同. 令 $\mathcal Z_k^*=\mathcal N(d_k^*)$, $\mathcal B_k^*=\mathcal R(d_{k+1}^*)$. 因此由 Theorem~\ref{thm:adjoint-kernel-range-orthogonality}, 有}
\[
\mathcal B^\perp=\mathcal Z^*,\quad\mathcal Z^\perp=\overline{\mathcal B}^*,\quad\mathcal B^{*\perp}=\mathcal Z,\quad\mathcal Z^{*\perp}=\overline{\mathcal B}.
\]

% block: dual-complex-chain-inclusion page:44 type:ordinary-paragraph
{从而}
\[
\mathcal B^*\subset\overline{\mathcal B}^*=\mathcal Z^\perp\subset\mathcal B^\perp=\mathcal Z^*,
\]
{因此}
\[
0\to V_n^*\xrightarrow{d_n^*}V_{n-1}^*\xrightarrow{d_{n-1}^*}\cdots\xrightarrow{d_1^*}V_0^*\to0
\]
{是链复形, 且 $(W,d^*)$ 是 Hilbert 复形. 若原 Hilbert 复形 $(W,d)$ 是闭的, 则由闭值域定理, 对偶复形 $(W,d^*)$ 也是闭的.}
